% Programming with InitechBase, Chapter 20: Order and Filters in Programs
% (AN-16).  Panels replayed (examples/p20-*.dot); sources cmd/set.c (SET
% ORDER, SET FILTER), cmd/query.c (GO, SKIP, BOF/EOF), the programs
% BYORDER.PRG and DEPTLIST.PRG.
\chapter{Order and Filters in Programs}
\Gates{test-pubs-examples, test-interp-set}
A program can change the order in which it meets the records, with SET
ORDER, and hide the records it does not want, with SET FILTER.  This chapter
uses both on the staff roster.

\begin{ChapterCovers}
\item Changing the master index in a program
\item Reading from the bottom
\item A filter from an answer
\item Leaving the settings as they were
\end{ChapterCovers}

\ATEHead{Preparing for This Chapter}
Start InitechBase with a fresh copy of the sample diskette.

\ATEHead{Two Orders}
\file{BYORDER.PRG} expects the roster opened with two indexes, on the name and
on the salary.  It reads the first three names in the order of the first, then
selects the second with SET ORDER and reads the three highest salaries from the
bottom with \cmd{SKIP -1}:
\ProgramListing{BYORDER.PRG}
\begin{Steps}
\item Type \type{USE STAFF INDEX STAFFNAM, STAFFSAL}.
\item Type \type{DO BYORDER}.
\end{Steps}
\DotExample{examples/p20-byorder}
After each SET ORDER the program moves with \cmd{GO TOP} or \cmd{GO BOTTOM}, and
before it ends it restores the order it found.

\ATEHead{A Filter from an Answer}
\file{DEPTLIST.PRG} lists the staff of the department the user names, through a
filter:
\ProgramListing{DEPTLIST.PRG}
\begin{Steps}
\item Type \type{USE STAFF}.
\item Type \type{DO DEPTLIST} and answer \type{acct}.
\item Type \type{DO DEPTLIST} and answer \type{mktg}.
\end{Steps}
\DotExample{examples/p20-deptlist}
When no record meets the filter, \cmd{GO TOP} leaves the pointer at the end of
the file, and \fn{EOF()} tells the program so.  The program removes the filter
before it ends.

\ATEHead{Where Do You Go From Here?}
Chapter~21 searches with \cmd{LOCATE} and \cmd{CONTINUE}.

\ATEHead{Summary}
\begin{enumerate}
\item SET ORDER TO \xp{n} changes the master index; move with \cmd{GO} after it.
\item \cmd{GO BOTTOM} and \cmd{SKIP -1} read an index from the end.
\item SET FILTER TO \xp{condition} hides records; SET FILTER TO removes the filter.
\item A program restores the settings it changed.
\end{enumerate}
